% generated by qpp/qpp2latex.py from qpp/examples/pairs.qpp; do not edit

The rationals in the kernel: pairs of masses up to the tie. Differences, squares, products of hypotheses, and the
words read by relation.

\section{Pairs}\label{sec:pairs}

\begin{lemma}[{squares of pairs}]\label{lem:pairsq}
A square of a difference, and the identity of one interval at order one.
\begin{stmt}
\sit{given}{x,y\in \Q}
\sit{(1)}{(x-y)^2+2 x y = x^2+y^2}
\sit{(2)}{s,b,l,r\in \Q:\ s^2\cdot (r^3-l^3)/3+s b\cdot (r^2-l^2)+b^2\cdot (r-l) = (r-l)/3\cdot ((s l+b)^2+(s l+b)\cdot (s r+b)+(s r+b)^2)}
\end{stmt}
\end{lemma}

\begin{proof}
Expansion.
\begin{plines}
\pline{(1)}{The goal follows by expanding.}
\pline{(2)}{The goal follows by expanding.}
\end{plines}
\end{proof}

\begin{lemma}[{bounds of pairs}]\label{lem:pairbound}
Bounds by a hypothesis times a square, and by a square alone.
\begin{stmt}
\sit{given}{l,r,P,Q\in \Q,\ l \le r}
\sit{(1)}{(r-l)/6\cdot (P^2+Q^2) \le (r-l)/3\cdot (P^2+P Q+Q^2)}
\sit{(2)}{P^2/2-Q^2 \le (P+Q)^2}
\end{stmt}
\end{lemma}

\begin{proof}
The gaps are $\frac{r-l}6(P+Q)^2$ and $\frac12(P+2Q)^2$.
\begin{plines}
\pline{(1)}{The goal follows from \(1/6\cdot (P+Q)^2\) times \((\mathrm{H}_{1})\).}
\pline{(2)}{The goal follows from \(1/2\cdot (P+2 Q)^2\).}
\end{plines}
\end{proof}

\begin{lemma}[{words read by relation}]\label{lem:pairwords}
The words read by relation: a positive part in a sum, a size and an indicator.
\begin{stmt}
\sit{given}{k\in \N,\ a_1,\dots,a_k,d_1,\dots,d_k\in \Q,\ x\in \Q}
\sit{(1)}{\textstyle\sum_{h=1}^k a_h \max(x-d_h,0) = \textstyle\sum_{h=1}^k a_h \mathbf 1[x > d_h]\cdot (x-d_h)}
\sit{(2)}{|x| \ge x}
\end{stmt}
\end{lemma}

\begin{proof}
Indicator: one case each way, the certificate in each.
\begin{plines}
\pline{(1)}{The goal follows term by term, by cases on the conditionals of the facts, then by computation in each case.}
\pline{(2)}{The goal follows by cases on the conditionals of the facts, then by computation in each case.}
\end{plines}
\end{proof}

\begin{lemma}[{larger and smaller}]\label{lem:pairmax}
The larger and the smaller of two pairs, an indicator and a size in two clauses, and a chain.
\begin{stmt}
\sit{given}{x,y,z\in \Q}
\sit{(1)}{\max(x,y) \ge x}
\sit{(2)}{x \le z,\ y \le z:\ \max(x,y) \le z}
\sit{(3)}{\max(x,y)+\min(x,y) = x+y}
\sit{(4)}{x \le y:\ \mathbf 1[x > y] = 0;\ x \ge 0:\ |x| = x}
\sit{(5)}{x \le y,\ y \le z:\ x \le y \le z}
\end{stmt}
\end{lemma}

\begin{proof}
The bounds of the words read by relation, and their cases.
\begin{plines}
\pline{(1)}{The goal follows from the facts at hand.}
\pline{(2)}{The goal follows by cases on the conditionals of the facts, then by computation in each case.}
\pline{(3)}{The goal follows by cases on the conditionals of the facts, then by computation in each case.}
\pline{(4)}{Clause 1: the goal follows by cases on the conditionals of the facts, then by computation in each case. Clause 2: the goal follows by cases on the conditionals of the facts, then by computation in each case.}
\pline{(5)}{Each part of the goal follows from the facts at hand.}
\end{plines}
\end{proof}

\begin{lemma}[{sums of pairs}]\label{lem:pairsums}
Sums term by term: positive parts, and a sum of squares at least 0.
\begin{stmt}
\sit{given}{k\in \N_{\ge1},\ a_1,\dots,a_k\in \Q_{\ge0},\ d_1,\dots,d_k,u_1,\dots,u_k\in \Q,\ x\in \Q}
\sit{(1)}{\textstyle\sum_{h=1}^k a_h \max(x-d_h,0) \ge 0}
\sit{(2)}{\textstyle\sum_{h=1}^k u_h^2 \ge 0}
\sit{(3)}{\textstyle\sum_{h=1}^k (\max(u_h,d_h)+\min(u_h,d_h)) = \textstyle\sum_{h=1}^k (u_h+d_h)}
\end{stmt}
\end{lemma}

\begin{proof}
Term by term.
\begin{plines}
\pline{(1)}{The goal follows term by term.}
\pline{(2)}{The goal follows term by term.}
\pline{(3)}{The goal follows term by term, by cases on the conditionals of the facts, then by computation in each case.}
\end{plines}
\end{proof}

\begin{lemma}[{sums by induction}]\label{lem:pairind}
Identities of sums by induction on the bound: linearity and a constant term.
\begin{stmt}
\sit{given}{k\in \N_{\ge1},\ m\in \N,\ u_1,\dots,u_k,v_1,\dots,v_k,w_1,\dots,w_m\in \Q,\ c\in \Q}
\sit{(1)}{\textstyle\sum_{h=1}^k (c u_h+v_h) = c\cdot (\textstyle\sum_{h=1}^k u_h)+(\textstyle\sum_{h=1}^k v_h)}
\sit{(2)}{\textstyle\sum_{h=1}^m (w_h-c) = \textstyle\sum_{h=1}^m w_h-m c}
\sit{(3)}{\textstyle\sum_{h=1}^k (u_h+c) = \textstyle\sum_{h=1}^k u_h+k c}
\end{stmt}
\end{lemma}

\begin{proof}
Each by induction on its bound.
\begin{plines}
\pline{(1)}{The goal follows by induction on the range of the sums.}
\pline{(2)}{The goal follows by induction on the range of the sums.}
\pline{(3)}{The goal follows by induction on the range of the sums.}
\end{plines}
\end{proof}

\begin{lemma}[{powers}]\label{lem:pairpow}
Powers with a variable exponent and a positive factor: the power-mean step of a heavy cluster.
\begin{stmt}
\sit{given}{N,P,q\in \Q_{\ge0},\ j\in \N_{\ge1}}
\sit{(1)}{q < 1,\ N^j \le q P^j:\ N \le P,\ (N = P\Rightarrow P = 0)}
\end{stmt}
\end{lemma}

\begin{proof}
Compare the powers; for the tie, cancel the factor $1-q$.
\begin{plines}
\pline{(1)}{The parts of the goal follow: (i) from the \(j\)-th powers; (ii) from the \(j\)-th powers, dividing by \(1-q\).}
\end{plines}
\end{proof}
